\section{Limitations and deployment boundary}
\label{sec:limitations}
The camera ladder crops and resamples existing orthographic imagery. It changes geographic footprint and assigned source pixels, but it cannot synthesize a new viewing angle, acquisition time, illumination condition, occluded surface, or detail below native resolution. Pre- and post-disaster channels undergo different resampling, and basemap context may come from a different date. The fixed flat-ground, nadir abstraction does not validate altitude-dependent image formation for a physical UAV.

The vehicle model uses simplified position updates without wind, attitude dynamics, acceleration limits, terrain and collision avoidance, communications, or battery discharge. Executed reobservations are measured software actions. State-derived horizontal and vertical displacements are geometric estimates. Of 3,360 final episodes, 42 end with a fly action lacking a subsequent recorded state, so complete-motion summaries are available for 3,318. Episode wall time is measured execution time on the experiment system, not measured flight time; render, ChangeOS, controller, and Qwen component latencies were not separately recorded. No PX4/Gazebo or physical-flight action mapping was performed. The study therefore concerns imagery-grounded active observation, not operational flight performance.

The ChangeOS checkpoint is fixed across policies but has prior exposure to related xBD data; the final policy comparison does not imply event-disjoint perception training. The final questions occupy 44 ROIs in only three disaster events. Although all 160 questions received the author's manual approval before online testing, this is not a two-annotator independent review or an agreement study, and 88 recorded ambiguity notices remain part of the audit trail. Answer frequencies and template families may still permit language-only shortcuts; the same-data development majority and question-only baselines cannot be transported as final-set estimates.

The generation protocol gives three seeded repeats, but active policies can make different numbers of later Qwen calls. Their paired outcomes thus compare entire decision pipelines rather than the isolated causal effect of one movement. The seven policies share a maximum reobservation cap, not equal executed action counts. No randomized cap sweep, full module-level timing, or complete layered GT/Qwen oracle suite is available. The interval for T1\_TASK versus A0\_HOLD includes both a small disadvantage and a useful advantage, and one of three event directions is negative. These boundaries preclude claiming stable superiority over holding or generalization to unseen disasters from the present sample.

Archived four-class allocation and VLN component results in Appendix~\ref{app:historical} use different populations and/or an earlier observation stack. They are historical diagnostics, not replications of the binary ChangeOS final experiment, and their calibration constants or performance numbers are not transferred into the main evidence chain.
